\chapter{Uso de Inteligência Artificial}\label{cap:uso-ia}

Este trabalho utilizou ferramentas de IA generativa em dois papéis distintos.
No primeiro, elas são objeto de estudo e instrumento experimental: o Cursor e o
OpenAI Codex com GPT-5.6 Luna compõem as condições C3--C6, e o Composer 2.5 Fast
atuou como avaliador da rubrica qualitativa. Esse uso integra a metodologia e
está descrito nos Capítulos~\ref{cap:metodologia} e~\ref{cap:desenvolvimento}.

No segundo papel, a IA foi empregada como ferramenta de apoio à produção do
trabalho. Foram utilizados, de forma intercalada, o GPT-5.6 Sol, por meio do
OpenAI Codex, e o Claude Sonnet 5, da Anthropic, nas
seguintes atividades:

\begin{itemize}
\item revisão ortográfica, gramatical e de clareza da linguagem do texto;
\item auxílio na construção e na correção dos \textit{scripts} de coleta,
normalização, avaliação e geração de gráficos e diagramas.
\end{itemize}

Os diagramas foram produzidos por um \textit{script} Python, construído com o
auxílio desses modelos, que descreve cada figura como código e gera arquivos
editáveis no Excalidraw. Não foram utilizados modelos de geração
de imagens: o conteúdo de cada figura foi definido e conferido pelo autor com
base no protocolo e nos dados do experimento.

Todo texto sugerido foi revisado e reescrito pelo autor, e os \textit{scripts}
foram validados pela execução e pela conferência de suas saídas com os dados
brutos preservados no repositório. A concepção do estudo, as decisões
metodológicas, a análise dos resultados e as conclusões são de
responsabilidade exclusiva do autor.
